\section{主板怎样从设计文件成为可上电实物}

原理图上的网络名称不会自动变成铜线，PCB 文件也不会自动变成一块能够安全上电的主板。从设计意图到实物至少经过电气设计、物理实现、制造资料冻结、裸板制造、元件装联、结构检查、电气测试和首轮上电八个阶段。每一阶段都要交付下一阶段能够执行和验收的资料；只把版图文件“发给工厂”并不能闭合工程责任。

\topic{原理图定义连接意图，约束定义这些连接怎样才算可用}

原理图给出器件、网络、电源域和功能关系，网表回答“哪些引脚应该相连”。BOM 进一步确定料号、封装、精度、耐压、温度等级和可替代关系。对 DDR、PCIe、时钟和开关电源等网络，还要给出差分阻抗、长度/时延匹配、参考平面、过孔数量、拓扑、电流密度与禁布区域；否则版图工具只知道连通，不知道怎样保持信号与电源完整性。

PCB stackup 决定各铜层、介质厚度和参考平面。布局先固定连接器、处理器、内存、VRM 和时钟等关键器件，再安排去耦、电源回路与高速通道。布线完成后要运行 DRC，并分别做信号完整性、电源完整性、热和可制造性检查。前文讲过的回流路径、阻抗突变、串扰和压降，在这里都变成具体的层、线宽、间距、过孔和铜面积。

\begin{pdfdiagram}[x=1cm,y=1cm]
  \node[hw state, minimum width=2.35cm] (schematic) at (1.25,4.8) {原理图/BOM\\网表与器件边界};
  \node[hw control, minimum width=2.35cm] (layout) at (4.1,4.8) {stackup/布局布线\\SI/PI/热约束};
  \node[hw state, minimum width=2.35cm] (release) at (6.95,4.8) {制造资料冻结\\版本与变更记录};
  \node[hw compute, minimum width=2.35cm] (pcb) at (9.8,4.8) {裸板制造\\成像、蚀刻、钻镀};

  \node[hw compute, minimum width=2.35cm] (assembly) at (9.8,2.4) {锡膏与贴装\\回流/选择焊};
  \node[hw control, minimum width=2.35cm] (inspect) at (6.95,2.4) {AOI/X-ray\\外观与隐藏焊点};
  \node[hw control, minimum width=2.35cm] (electrical) at (4.1,2.4) {ICT/飞针/JTAG\\开短路与连接};
  \node[hw state, minimum width=2.35cm] (unpowered) at (1.25,2.4) {不上电检查\\阻值、极性、短路};

  \node[hw control, minimum width=2.35cm] (limited) at (1.25,0.0) {限流首轮上电\\先看待机与主轨};
  \node[hw state, minimum width=2.35cm] (clockreset) at (4.1,0.0) {时钟/复位/SPI\\逐级寻找证据};
  \node[hw compute, minimum width=2.35cm] (boot) at (6.95,0.0) {固件与自检\\记录里程碑};
  \node[hw state, minimum width=2.35cm] (releaseprod) at (9.8,0.0) {工程样机放行\\问题回写设计};

  \draw[hw bus] (schematic) -- (layout);
  \draw[hw bus] (layout) -- (release);
  \draw[hw bus] (release) -- (pcb);
  \draw[hw bus] (pcb) -- (assembly);
  \draw[hw bus] (assembly) -- (inspect);
  \draw[hw bus] (inspect) -- (electrical);
  \draw[hw bus] (electrical) -- (unpowered);
  \draw[hw bus] (unpowered) -- (limited);
  \draw[hw bus] (limited) -- (clockreset);
  \draw[hw bus] (clockreset) -- (boot);
  \draw[hw bus] (boot) -- (releaseprod);
\end{pdfdiagram}
\diagramnote{从设计意图到工程样机的三层闭环。上排把原理图和物理约束冻结成可制造资料，中排完成装联并用互补测试寻找结构与连接缺陷，下排在限流条件下逐级建立电气、时钟、复位和固件证据。任何制造问题都要回写到明确的设计版本，不能只在样板上临时补线后遗忘。}

\topic{制造包是一组相互核对的合同}

常见制造资料包括铜层与阻焊/字符层图形、NC drill 钻孔文件、板框与机械尺寸、stackup 和受控阻抗说明；装联资料还包括坐标文件、钢网开口、装配图、BOM 和特殊工艺说明。Gerber 与 ODB++ 等格式只是承载方式，关键是所有文件必须来自同一次受控 release，并有可追溯版本。

\begin{longtable}{L{0.22\linewidth}L{0.29\linewidth}L{0.39\linewidth}}
\toprule
交付物 & 主要回答的问题 & 单独拿它仍不能证明 \\
\midrule
原理图与网表 & 哪些引脚应连接，器件承担什么功能 & 实际走线满足阻抗、回流和电流要求 \\\tablerowrule
PCB 数据与 stackup & 每层图形、板框、孔和材料怎样实现 & 装了正确元件，也不能证明焊接合格 \\\tablerowrule
BOM 与替代规则 & 料号、封装、参数和允许替代是什么 & 贴片机坐标、方向和工艺温度正确 \\\tablerowrule
坐标/装配图/钢网资料 & 元件放在哪里、怎样定向、焊膏怎样印刷 & 隐藏 BGA 焊点和内部过孔可靠 \\\tablerowrule
测试点与边界扫描描述 & 哪些网络可被电气激励和观察 & 高速眼图、满载电源与固件功能全部合格 \\\tablerowrule
release 清单与散列值 & 本批生产使用哪一套受控文件 & 设计本身没有逻辑错误 \\
\bottomrule
\end{longtable}

正式下单前要做 DFM 与 DFT 评审。DFM 检查最小线宽间距、孔径、环宽、阻焊桥、拼板、器件边缘距离和装联可达性；DFT 检查关键电源、复位、时钟、总线和调试链是否有可观察点。没有测试点的网络不一定错误，但出现故障时会显著扩大定位成本。

\topic{裸板制造先形成互连，再保护与成形}

多层板从各内层图形的成像与蚀刻开始，经过层压形成整体；钻孔后对孔壁沉铜和电镀，使不同铜层通过过孔连通。外层图形继续成像、蚀刻，再施加阻焊、表面处理和字符。最后铣出板框并做裸板电测。具体先后会随 HDI、盲埋孔、背钻和材料体系变化，但“形成图形—建立层间连接—保护表面—电测连通”的逻辑不变。

受控阻抗不能只靠最终线宽猜测。板厂要根据实际介质厚度、铜厚与材料介电特性调整，并用阻抗 coupon 或过程测试核对。设计方接受板厂补偿时必须同步记录，否则回来的板可能通过制造 DRC，却偏离最初的时延与阻抗预算。

\topic{装联把正确器件放到正确方向，并形成可重复焊点}

表面贴装通常先印刷锡膏、用贴片机放置元件，再经过受控温度曲线回流。通孔器件可能使用选择焊、波峰焊或手工补焊。回流温度曲线要兼顾焊膏活化、器件耐温、板翘和空洞；大型 BGA、连接器和厚铜区域的热容量不同，不能只看炉温设定值。

AOI 适合检查可见元件的缺件、偏移、极性和焊点外观；BGA 焊球位于封装下方，通常需要 X-ray 观察桥连、开路迹象和空洞分布。显微镜、切片与染色分析可在疑难失效中给出更直接证据，但往往具有破坏性，应保留代表样品和问题批次信息。

\begin{longtable}{L{0.17\linewidth}L{0.27\linewidth}L{0.27\linewidth}L{0.19\linewidth}}
\toprule
检查手段 & 擅长发现 & 不擅长单独证明 & 合适阶段 \\
\midrule
AOI & 缺件、偏移、极性、可见桥连与焊量异常 & BGA 下隐藏焊点、内部层间开路 & 回流后快速筛查 \\\tablerowrule
X-ray & BGA 桥连、明显开路迹象、空洞与异物 & 焊点长期机械寿命、所有细微裂纹 & 隐藏焊点检查 \\\tablerowrule
ICT/飞针 & 电源短路、网络开路、阻容值和部分器件方向 & 高速运行眼图与完整固件功能 & 上电前电气筛查 \\\tablerowrule
JTAG 边界扫描 & 可控数字引脚间的连通、焊接开短路 & 不在扫描链中的模拟和高速内部功能 & 装联后、功能启动前 \\\tablerowrule
功能测试 & 真实时钟、协议、负载和软件路径 & 未覆盖条件下的潜在制造缺陷 & 首轮上电后 \\
\bottomrule
\end{longtable}

\topic{首轮上电以前，先在无能量条件下排除明显风险}

装联板不应拿到手就接额定大功率电源。先核对板号与 release、关键器件方向、跳线与电源配置，检查是否残留锡珠和金属异物；再在断电状态测量各电源轨对地阻值，并与良板、仿真或设计预期比较。低压大电流核心轨本来可能呈现较低阻值，不能仅凭蜂鸣档判定短路，应结合电路结构、注入电压限制和热像证据。

还要确认输入保护、保险或电子熔丝、反接保护和各稳压器 enable 的默认状态。可拆分的电源域应分阶段接入；外部烧录器、示波器地线和串口适配器也可能通过信号脚反向供电，连接前必须检查地电位和 I/O 耐压。

\topic{限流上电的目标是获得第一条可靠证据}

实验电源的电压和限流值应根据本板预计的待机电流、启动浪涌、电容和电源路径预算设置，而不是套用一个固定数字。限流过高失去保护作用，过低则可能让正常电源反复打嗝，产生“板坏了”的假象。首次可只使能待机域，观察输入电流和待机轨；确认无异常发热后，再允许主电源和各负载点按既定顺序启动。

\begin{longtable}{L{0.09\linewidth}L{0.24\linewidth}L{0.29\linewidth}L{0.28\linewidth}}
\toprule
阶段 & 允许注入的能量 & 要建立的证据 & 异常时停止在哪里 \\
\midrule
B0 & 完全断电，仅用受限小信号测量 & 版本、极性和关键阻值合理 & 不连接主电源 \\\tablerowrule
B1 & 限流输入，仅待机域 & 输入电流受控，待机轨与 EC/SIO 时钟存在 & 关断并找短路、反装或错误料 \\\tablerowrule
B2 & 允许主电源路径启动 & 各轨爬升、PGOOD 和时序合格 & 在首个异常电源域停止后级 enable \\\tablerowrule
B3 & 开放参考时钟与复位树 & 晶振、PLL lock 和复位释放顺序成立 & 不让 CPU/高速器件进入未知执行 \\\tablerowrule
B4 & 允许 Boot ROM/SPI 访问 & 复位向量、片选、时钟和数据活动存在 & 保存 trace，不盲目反复重启 \\\tablerowrule
B5 & 运行最小固件自检 & POST、DRAM 训练和枚举逐步闭合 & 记录最后成功阶段与第一失败条件 \\\tablerowrule
B6 & 逐步增加设备和负载 & 满载电源、温升、高速接口与长稳合格 & 问题回写到设计、工艺或测试覆盖 \\
\bottomrule
\end{longtable}

\topic{能启动一次不等于制造与设计已经验收}

一块板成功打印串口，只证明本次样品在当前条件下走通了某条功能路径。工程样机还要核对电源纹波与瞬态、时钟抖动、高速链路误码、DRAM 压力、热稳态、复位/掉电、边界电压温度和重复启动。制造批次还要看良率分布，而不是只分析一块“黄金样板”。

故障要先归到证据层：原理图连接错误、版图/约束错误、裸板开短路、错料反装、焊接缺陷、固件配置错误或器件本身失效。飞线、换料和手工补焊可以验证假设，但最终修复必须进入原理图、PCB、BOM、工艺文件或测试程序的新 release。这样下一批板继承的是被修正的设计，而不是只有某位工程师记得的一次现场操作。

到这里，主板才从“设计文件中的连接关系”变成“具有可追溯制造证据、能够受控上电并运行固件的物理平台”。后文“从电源键到复位向量”的 trace，可以建立在这块已经通过无源检查与分级上电的板上；随后讲到的多核启动，则继续把一颗已取指的处理器扩展成可被操作系统完整使用的计算平台。
